\chapter{Lorentzian Green Functions and Analytic Continuation}
\label{ch:lorentzian-green-functions}
\ConventionsMixed


\section{Time-Ordered Boundary Values}

Let \(x^\mu=(x^0,\vec x)\in\mathbb R^{1,D-1}\) and use the mostly-plus
metric
\[
  \eta_{\mu\nu}=\operatorname{diag}(-,+,\ldots,+),
  \qquad
  k^2=\eta_{\mu\nu}k^\mu k^\nu.
\]
Write \(M=\mathbb R^{1,D-1}\).  The framework in this chapter is a scalar
Wightman field presentation in a vacuum sector: a Hilbert space \(\Hilb\), a
vacuum vector \(\vac\), a strongly continuous unitary representation of
translations with self-adjoint generators \(P^\mu\) whose joint spectrum lies
in the closed forward cone
\[
  \overline V_+=\{p\in M:p^0\ge0,\ p^2\le0\},
\]
and a scalar
operator-valued distribution \(\widehat\phi\) defined on a common invariant
dense domain.  Point-field notation denotes the kernel notation for smeared
operator-valued distributions.

For \(r\ge1\), \(\mathcal S(M^r)\) denotes the Schwartz space on \(M^r\), and
\(\mathcal S'(M^r)\) its continuous dual.  This chapter works in the tempered
distribution category:
\[
  W_r,\ G_r\in\mathcal S'(M^r).
\]
Every compactly supported smearing used to define a local field matrix element
is also a Schwartz test function, so this convention is compatible with the
local operator-valued-distribution notation.  After translation invariance is
used, analytic boundary values are paired with test functions in the relative
variables.  Thus a holomorphic function \(F\) on a tube over \(M^{r-1}\) has
boundary value \(T\in\mathcal S'(M^{r-1})\) when, for every
\(f\in\mathcal S(M^{r-1})\),
\[
  \lim_{y\to0}
  \int_{M^{r-1}}\dd^{D(r-1)}\xi\,
  F(\xi-\ii y)f(\xi)
  =
  \langle T,f\rangle ,
\]
with \(y\) approaching the real edge inside the specified tube cone.

The Wightman \(n\)-point distribution is
\[
  W_n(x_1,\ldots,x_n)
  =
  \langle \vac|
  \widehat\phi(x_1)\cdots\widehat\phi(x_n)
  |\vac\rangle
  \in \mathcal S'(M^n).
\]
For a permutation \(\pi\in S_n\), set
\[
  M_\pi^n
  =
  \{(x_1,\ldots,x_n):x_{\pi(1)}^0>\cdots>x_{\pi(n)}^0\}.
\]
A vacuum time-ordered Green function for this field is a distribution
\[
  G_n\in\mathcal S'(M^n)
\]
whose restriction to \(M_\pi^n\) is
\[
  G_n(x_1,\ldots,x_n)\big|_{M_\pi^n}
  =
  \langle \vac|
  \widehat\phi(x_{\pi(1)})\cdots
  \widehat\phi(x_{\pi(n)})
  |\vac\rangle .
\]
Equivalently, on \(M_\pi^n\) away from coincident insertion points,
\[
  G_n(x_1,\ldots,x_n)
  =
  \langle \vac|
  \widehat\phi(x_{\pi(1)})\cdots
  \widehat\phi(x_{\pi(n)})
  |\vac\rangle .
\]
Local commutativity makes the restrictions compatible across spacelike
equal-time boundaries.  A choice of \(G_n\) is therefore an extension across
partial diagonals \(x_{a_1}=\cdots=x_{a_r}\).  The possible differences
between two such extensions are local distributions supported on these
diagonals.  In perturbation theory this extension datum is fixed together
with the renormalization prescription for local products.  The compact
operator notation
\[
  G_n(x_1,\ldots,x_n)
  =
  \langle \vac|
  \operatorname T\{
    \widehat\phi(x_1)\cdots\widehat\phi(x_n)
  \}
  |\vac\rangle
\]
refers to this specified distribution.

For the free scalar field of mass \(m>0\), the two-point time-ordered Green
function has Fourier representation
\[
  \Delta_F(x;m^2)
  =
  -\ii
  \int \frac{\dd^D k}{(2\pi)^D}\,
  \frac{\ee^{\ii k\cdot x}}{k^2+m^2-\ii\epsilon},
  \qquad
  \epsilon>0.
\]
At fixed \(\vec k\), with
\[
  \omega_{\vec k}=\sqrt{\vec k^{\,2}+m^2},
\]
the poles in the complex \(k^0\)-plane are
\[
  k^0=-\omega_{\vec k}+\ii0,
  \qquad
  k^0=+\omega_{\vec k}-\ii0 .
\]
This is the Feynman boundary value of the meromorphic function of \(k^0\).

\section{Tube Analyticity From the Spectrum Condition}

The \(i\epsilon\) prescription is the two-point perturbative expression of a
more general analytic fact.  Let
\[
  W_n(x_1,\ldots,x_n)
  =
  \langle \vac|
  \widehat\phi(x_1)\cdots\widehat\phi(x_n)
  |\vac\rangle
\]
be a Wightman distribution.  By translation invariance write it in the
relative variables
\[
  \xi_j=x_j-x_{j+1},
  \qquad j=1,\ldots,n-1.
\]
Let
\[
  w_n(\xi_1,\ldots,\xi_{n-1})\in\mathcal S'(M^{n-1})
\]
be the corresponding relative-coordinate distribution.  Its Fourier transform
\(\widetilde w_n\) is defined with the convention
\[
  w_n(\xi)
  =
  \int\frac{\dd^{D(n-1)}p}{(2\pi)^{D(n-1)}}\,
  \ee^{\ii\sum_j p_j\cdot\xi_j}\widetilde w_n(p)
\]
in distribution notation.  Equivalently, for \(f\in\mathcal S(M^{n-1})\),
\[
  \langle w_n,f\rangle
  =
  \left\langle \widetilde w_n,\check f\right\rangle,
  \qquad
  \check f(p)
  =
  \int\frac{\dd^{D(n-1)}\xi}{(2\pi)^{D(n-1)}}\,
  \ee^{\ii\sum_jp_j\cdot\xi_j}f(\xi).
\]
The spectrum condition gives the support inclusion
\[
  \operatorname{supp}\widetilde w_n
  \subset
  \overline V_+^{\,n-1}.
\]
For \(\eta=(\eta_1,\ldots,\eta_{n-1})\in V_+^{n-1}\), the function
\(\exp(\sum_j p_j\cdot\eta_j)\) decreases exponentially on every closed
subcone of \(\overline V_+^{\,n-1}\) which contains the support of
\(\widetilde w_n\).  Multiplication of \(\widetilde w_n\) by this function is
therefore a well-defined tempered distribution after inserting any smooth
cutoff equal to one on a conic neighborhood of the support; the result is
independent of the cutoff.  The inverse Fourier transform is a holomorphic
function of
\begin{equation}
  z_j=\xi_j-\ii\eta_j,
  \qquad
  \eta_j\in V_+,\qquad j=1,\ldots,n-1,
  \label{eq:wightman-forward-tube}
\end{equation}
where \(V_+\) is the open future light cone.  In distribution notation this
Fourier-Laplace transform is written as
\[
  \mathcal W_n(z)
  =
  \frac{1}{(2\pi)^{D(n-1)}}
  \left\langle
    \widetilde w_n(p),
    \exp\!\left\{\ii\sum_{j=1}^{n-1}p_j\cdot z_j\right\}
  \right\rangle .
\]
Indeed, for
\(p\in\overline V_+\) and \(\eta\in V_+\), the mostly-plus inner product
satisfies
\[
  p\cdot\eta\le0,
\]
so the factor
\[
  \exp\{\ii p\cdot(\xi-\ii\eta)\}
  =
  \exp\{\ii p\cdot\xi+p\cdot\eta\}
\]
is damped rather than enlarged.
The boundary value is \(w_n\) in \(\mathcal S'(M^{n-1})\):
\[
  \lim_{\eta\to0,\ \eta_j\in V_+}
  \int_{M^{n-1}}\dd^{D(n-1)}\xi\,
  \mathcal W_n(\xi-\ii\eta)f(\xi)
  =
  \langle w_n,f\rangle,
  \qquad f\in\mathcal S(M^{n-1}).
\]

Different operator orderings give holomorphic functions in corresponding
permuted tubes.  Local commutativity identifies their boundary values on real
regions where the exchanged neighboring points are spacelike separated.

\begin{theorem}[Distributional edge-of-the-wedge statement used here]
\label{thm:distributional-edge-of-the-wedge}
Let \(U\subset\mathbb R^N\) be open, and let \(C_+\) and \(C_-\) be open
convex cones in \(\mathbb R^N\).  Suppose \(F_\pm\) are holomorphic on the
tubes \(U+\ii C_\pm\).  Assume the following polynomial boundary growth: for
every compact \(K\subset U\) and every closed cone \(C'_\pm\) whose nonzero
vectors lie in \(C_\pm\), there are constants \(A,N<\infty\) such that
\[
  |F_\pm(x+\ii y)|
  \le
  A(1+|x|)^N |y|^{-N},
  \qquad
  x\in K,\quad y\in C'_\pm,\quad 0<|y|<1.
\]
These estimates define distributional boundary values
\[
  b_\pm\in\mathcal D'(U).
\]
If \(b_+=b_-\) on a nonempty open real subset \(O\subset U\), then there is a
connected complex neighborhood \(\mathcal N\) of \(O\) and a holomorphic
function \(F\) on \(\mathcal N\) whose restrictions agree with \(F_+\) and
\(F_-\) wherever the original tubes meet \(\mathcal N\).
\end{theorem}

In the Wightman application, the polynomial growth hypothesis comes from the
tempered Fourier-Laplace construction above, and the equality of boundary
values comes from local commutativity on the spacelike real edge.  The theorem
therefore supplies the common holomorphic continuation of the permuted tube
functions across those spacelike real regions.  Time-ordered Green functions
are boundary values obtained by approaching real Lorentzian points from
imaginary time orderings compatible with the chosen operator order.  For two
points, this general construction is exactly the pole displacement
\[
  \frac1{k^2+m^2-\ii\epsilon}
\]
appearing in the Feynman propagator.

\section{Boundary-Value Dictionary}

The preceding tube statement gives a precise dictionary among Wightman,
time-ordered, and Euclidean functions.  Fix an ordering \(\pi\) and real
Lorentzian points with
\[
  x_{\pi(1)}^0>x_{\pi(2)}^0>\cdots>x_{\pi(n)}^0 .
\]
Choose positive numbers
\[
  \epsilon_1>\epsilon_2>\cdots>\epsilon_n>0
\]
and set
\[
  z_{\pi(j)}=x_{\pi(j)}-\ii\epsilon_j e_0,
  \qquad e_0=(1,\vec0).
\]
Then
\[
  z_{\pi(j)}-z_{\pi(j+1)}
  =
  x_{\pi(j)}-x_{\pi(j+1)}
  -
  \ii(\epsilon_j-\epsilon_{j+1})e_0
\]
lies in the forward tube for the \(j\)-th relative variable.  The
time-ordered boundary value on the region \(M_\pi^n\) is
\[
  G_n(x_1,\ldots,x_n)\big|_{M_\pi^n}
  =
  \lim_{\epsilon_1>\cdots>\epsilon_n\downarrow0}
  \mathcal W_{n,\pi}(z_{\pi(1)},\ldots,z_{\pi(n)}),
\]
where \(\mathcal W_{n,\pi}\) denotes the holomorphic continuation of the
Wightman ordering displayed by \(\pi\).  The limit is a distributional
boundary value: it is taken after pairing with a Schwartz test function in
the relative variables.

Euclidean Schwinger functions arise from the same holomorphic functions by
restriction to imaginary time.  For Euclidean points
\[
  x_{E,a}=(\tau_a,\vec x_a),
\]
with
\[
  \tau_{\pi(1)}>\tau_{\pi(2)}>\cdots>\tau_{\pi(n)},
\]
define
\[
  z_{\pi(j)}^0=-\ii\tau_{\pi(j)},
  \qquad
  \vec z_{\pi(j)}=\vec x_{\pi(j)} .
\]
Then
\[
\begin{aligned}
  &S_n(x_{E,1},\ldots,x_{E,n})
    \big|_{\tau_{\pi(1)}>\cdots>\tau_{\pi(n)}}  \\
  &\qquad =
  \mathcal W_{n,\pi}(z_{\pi(1)},\ldots,z_{\pi(n)}).
\end{aligned}
\]
Conversely, an OS-admissible Schwinger hierarchy reconstructs Wightman
boundary values by the positive-time quotient and analytic-continuation
theorem of Chapter~\ref{ch:volume-iv-osterwalder-schrader-reconstruction}.
Thus the boundary-value dictionary has two certified entrances: spectral
Wightman data with tube analyticity, and Euclidean data satisfying the OS
positivity, covariance, regularity, and growth hypotheses.

For \(n=2\), the dictionary gives the Feynman prescription.  If
\(x^0>0\), the ordered boundary value is obtained with the positive-energy
pole below the real \(k^0\)-axis and the negative-energy pole above it:
\[
  k^0=+\omega_{\vec k}-\ii0,
  \qquad
  k^0=-\omega_{\vec k}+\ii0.
\]
If \(x^0<0\), the contour closure selects the opposite residue in the same
single distribution \(\Delta_F\).  This two-point pole rule is the momentum
space encoding of the ordered tube boundary value above.

\section{Lorentzian Feynman Rules for Green Functions}

Consider the same scalar \(\phi^4\) interaction as in the Euclidean chapter,
with Lorentzian interaction density
\[
  \mathcal L_{\mathrm{int}}
  =
  -\frac{g}{4!}\phi^4 .
\]
The perturbative rules below compute time-ordered Green functions. They use
Lorentzian momenta on every line.

\begin{center}
\begin{tikzpicture}[scale=0.9,>=Stealth]
  \draw[purple!80!black,very thick,->] (-1.25,0)--(1.25,0)
    node[midway,above] {$k$};
  \node at (2.0,0) {$=$};
  \node at (3.35,0)
    {$\displaystyle \frac{-\ii}{k^2+m^2-\ii\epsilon}$};

  \begin{scope}[xshift=7.2cm]
    \fill[blue!80!black] (0,0) circle (2.5pt);
    \foreach \ang/\lab in {135/k_1,45/k_2,-45/k_3,-135/k_4} {
      \draw[purple!80!black,very thick,->]
        ({1.0*cos(\ang)},{1.0*sin(\ang)}) -- (0,0);
      \node at ({1.34*cos(\ang)},{1.34*sin(\ang)}) {$\lab$};
    }
    \node at (1.85,0) {$=$};
    \node[align=left,anchor=west] at (2.25,0)
      {$\displaystyle -\ii g\,(2\pi)^D
        \delta^{(D)}(k_1+k_2+k_3+k_4)$};
  \end{scope}
\end{tikzpicture}
\end{center}

This is the labelled-leg convention matching the Euclidean convention of the
previous chapter. Direct expansion of the interaction density instead assigns
\[
  -\frac{\ii g}{4!}\,(2\pi)^D
  \delta^{(D)}(k_1+k_2+k_3+k_4)
\]
to an unlabelled four-field insertion and recovers the labelled rule by
counting contractions.
In the labelled-leg convention used below, the factor \(4!\) associated with
assigning the four incident half-edges to a quartic vertex has already been
absorbed into the vertex rule \(-\ii g\).  Therefore a connected graph
\(\Gamma\) with labelled external legs carries only the residual factor
\[
  \frac{1}{|\operatorname{Aut}_{\mathrm{ext}}\Gamma|},
\]
where \(\operatorname{Aut}_{\mathrm{ext}}\Gamma\) is the automorphism group of
the graph that fixes every external label and preserves the incidence
relations.  For the tadpole graph below this residual group has order \(2\),
coming from the interchange of the two internal half-edges forming the loop;
this is why the tadpole coefficient is \(-\ii g/2\), not \(-\ii g\).

\section{Loop Rotation and the Tadpole}

The one-loop tadpole contribution to the Lorentzian two-point self-energy
insertion is
\[
\begin{tikzpicture}[baseline=-0.45ex,scale=0.75,>=Stealth]
  \draw[purple!80!black,very thick,->] (-1.1,0)--(0,0)
    node[midway,below] {$k$};
  \draw[purple!80!black,very thick,->] (1.1,0)--(0,0)
    node[midway,below] {$k$};
  \fill (0,0) circle (2.2pt);
  \draw[purple!80!black,very thick,->]
    (0,0) .. controls (-0.35,0.9) and (0.35,0.9) .. (0,0)
    node[midway,above] {$\ell$};
\end{tikzpicture}
\quad=\quad
-\frac{\ii g}{2}
\int \frac{\dd^D\ell}{(2\pi)^D}\,
\frac{-\ii}{\ell^2+m^2-\ii\epsilon}.
\]
The factor \(1/2\) is the same tadpole symmetry factor as in Euclidean
signature.  This is a combinatorial statement: the Wick contraction count
depends only on the polynomial \(\phi^4\) and on the residual graph
automorphism, while the signature-dependent factors enter separately through
the labelled Lorentzian vertex \(-\ii g\) and propagator numerator \(-\ii\).
The loop energy \(\ell^0\) is integrated along the real axis with
the Feynman pole displacement. The Euclidean integral is obtained by rotating
the contour to \(\ell^0=\ii \ell_E^D\), with the contour passing between the
two displaced poles.

\begin{center}
\begin{tikzpicture}[scale=0.95,>=Stealth]
  \draw[->,line width=0.65pt] (-3.1,0)--(3.1,0)
    node[right] {$\operatorname{Re}\ell^0$};
  \draw[->,line width=0.65pt] (0,-2.05)--(0,2.05)
    node[above] {$\operatorname{Im}\ell^0$};
  \draw[blue!70!black,very thick,->] (-2.75,0)--(2.75,0);
  \draw[green!50!black,very thick,->] (0,-1.75)--(0,1.75);
  \node[green!50!black,anchor=west,font=\small] at (0.18,1.82)
    {$\ell^0=\ii\ell_E^D$};
  \node[red!80!black] at (-1.35,0.35) {$\times$};
  \node[red!80!black] at (1.35,-0.35) {$\times$};
  \node[red!80!black,anchor=east] at (-1.52,0.92)
    {$-\omega_{\vec\ell}+\ii0$};
  \node[red!80!black,anchor=west] at (1.52,-1.02)
    {$+\omega_{\vec\ell}-\ii0$};
  \draw[orange!85!black,very thick,->]
    (0.38,1.45) .. controls (1.25,1.12) and (1.95,0.65) .. (2.45,0.08);
  \draw[orange!85!black,very thick,->]
    (-0.38,-1.45) .. controls (-1.25,-1.12) and (-1.95,-0.65) .. (-2.45,-0.08);
  \node[orange!85!black,font=\small,anchor=west] at (2.25,1.35)
    {contour rotation};
\end{tikzpicture}
\end{center}

With this orientation,
\[
\begin{gathered}
  \ell^0=\ii\ell_E^D,\qquad
  \dd\ell^0=\ii\,\dd\ell_E^D,\qquad
  \dd^D\ell=\ii\,\dd^D\ell_E,\\
  \ell^2+m^2-\ii\epsilon
  =
  -(\ell^0)^2+\vec\ell^{\,2}+m^2-\ii\epsilon
  \longrightarrow
  \ell_E^2+m^2 .
\end{gathered}
\]
Therefore the measure and propagator numerator together give
\[
  \dd^D\ell\,\frac{-\ii}{\ell^2+m^2-\ii\epsilon}
  \longrightarrow
  \ii\,\dd^D\ell_E\,\frac{-\ii}{\ell_E^2+m^2}
  =
  \dd^D\ell_E\,\frac{1}{\ell_E^2+m^2}.
\]
Including the vertex factor, this gives
\[
\begin{aligned}
  -\frac{\ii g}{2}
  \int \frac{\dd^D\ell}{(2\pi)^D}
  \frac{-\ii}{\ell^2+m^2-\ii\epsilon}
  &=
  \ii\left[
    -\frac{g}{2}
    \int \frac{\dd^D\ell_E}{(2\pi)^D}
    \frac{1}{\ell_E^2+m^2}
  \right].
\end{aligned}
\]
Thus the Lorentzian two-point insertion is \(\ii\Sigma(k)\), where
\(\Sigma(k)\) is the first-sheet analytic continuation of the Euclidean
self-energy \(\Sigma_E(k_E)\) with the sign convention of
Chapter~\ref{ch:perturbative-green-functions}.
More explicitly, in the same regulator and the same local two-point
subtraction convention, the denominator self-energy obeys
\begin{equation}
  \Sigma\bigl(k^0=\ii k_E^D,\vec k=\vec k_E\bigr)
  =
  \Sigma_E(k_E),
  \label{eq:lorentzian-euclidean-self-energy-match}
\end{equation}
where \(\Sigma_E\) is defined by
Eq.~\eqref{eq:euclidean-self-energy}.  The factor \(\ii\) displayed above
belongs to the Lorentzian two-point insertion in the Dyson series; it is not an
extra sign in the scalar function \(\Sigma\) appearing in the denominator of
\(\widetilde G(k)\).

The contour rotation just used is a special case with no external energy
flowing through the loop.  The same direct rotation is also legitimate for
loop integrals with Euclidean external momenta, where the external energies
lie on the imaginary axis and the Feynman poles can be separated continuously
from the rotating contours.  For general Lorentzian timelike external momenta,
one must instead analyze the full contour deformation in the complex loop
energy variables.  The obstruction is a pinch singularity: poles from different
propagators can trap the integration cycle so that it cannot be deformed to
the Euclidean cycle without crossing singularities.

For example, the Lorentzian continuation of the order-\(g^2\) sunset
self-energy contains the denominator product
\[
  \frac{1}{
  (q_1^2+m^2-\ii0)(q_2^2+m^2-\ii0)
  \bigl((k-q_1-q_2)^2+m^2-\ii0\bigr)} .
\]
At \(\vec k=0\), the Landau equations have a solution with all three internal
lines on shell and carrying positive energy when
\[
  (k^0)^2\ge (3m)^2 .
\]
At and above this three-particle threshold the \(q_1^0,q_2^0\) contours are
pinched.  Thus the Euclidean sunset integral defines a first-sheet analytic
function away from its Landau singularities, and the physical Lorentzian value
above threshold is the specified boundary value on the cut, not the result of
a naive independent rotation of each loop energy.  The discontinuity across
such cuts is computed by the Cutkosky rules, which are the perturbative
expression of the same spectral and unitarity structure developed in
Chapter~\ref{ch:volume-ii-analyticity-crossing-landau}.

\section{Full Two-Point Function and First Sheet}

The perturbative time-ordered two-point function is written as
\[
  G(x,y)
  =
  \int\frac{\dd^D k}{(2\pi)^D}\,
  \ee^{\ii k\cdot(x-y)}
  \widetilde G(k),
\]
where
\[
  \widetilde G(k)
  =
  \frac{-\ii}{k^2+m^2-\ii\epsilon-\Sigma(k)} .
\]
The function \(\Sigma(k)\) is initially computed for Euclidean external
momenta, where loop integrals have Euclidean denominators. Using
\[
  k^0=\ii k_E^D ,
\]
the Euclidean external energy lies on the imaginary \(k^0\)-axis. The
matching identity \eqref{eq:lorentzian-euclidean-self-energy-match} gives
\[
  \widetilde G(\ii k_E^D,\vec k_E)
  =
  -\ii\,\widetilde G_E(k_E),
\]
because \(k^2=k_E^2\) at \(k^0=\ii k_E^D\) with the mostly-plus convention.
The Lorentzian value is the boundary value obtained by analytic continuation of
\(k^0\) from this axis to the prescribed side of the real axis without crossing
the first-sheet cuts.

The spectral representation fixes the location of the first singular
thresholds. Suppose the continuum part \(d\rho_{\mathrm c}\) in the two-point
spectral measure has lowest invariant mass \(M\), meaning
\[
  M^2
  =
  \inf \operatorname{supp} d\rho_{\mathrm c}
\]
in the K\"all\'en--Lehmann variable \(s=\mu^2\).  Under the two-particle
threshold assumption stated in Chapter~\ref{ch:kallen-lehmann}, namely
\(\operatorname{supp}d\rho_{\mathrm c}\subset[4m^2,\infty)\) with the
\(\phi\phi\) continuum saturating the lower endpoint, this notation gives
\[
  M^2=4m^2,\qquad M=2m .
\]
If the first continuum channel has a different threshold, \(M\) denotes that
actual square-root threshold. At fixed spatial momentum \(\vec k\), define
\[
  E_M(\vec k)=\sqrt{\vec k^{\,2}+M^2}.
\]
The continuum denominators have singular support when
\[
  (k^0)^2=\vec k^{\,2}+s,
  \qquad s\ge M^2.
\]
Consequently the first sheet has branch cuts beginning at
\[
  k^0=+E_M(\vec k),
  \qquad
  k^0=-E_M(\vec k).
\]
The Feynman boundary value on the positive-energy cut is reached from the
upper half-plane. Thus the value of \(\Sigma(k)\) at real positive \(k^0\) is
the boundary value approached from \(\operatorname{Im}k^0>0\) on the first
sheet. The negative-energy cut is reached with the conjugate placement.

\begin{center}
\begin{tikzpicture}[scale=0.95,>=Stealth]
  \draw[->,line width=0.65pt] (-3.15,0)--(3.15,0)
    node[right] {$\operatorname{Re}k^0$};
  \draw[->,line width=0.65pt] (0,-2.05)--(0,2.05)
    node[above] {$\operatorname{Im}k^0$};
  \draw[red!80!black,very thick,decorate,
    decoration={zigzag,segment length=4pt,amplitude=1.2pt}]
    (-3.0,0)--(-1.05,0);
  \draw[red!80!black,very thick,decorate,
    decoration={zigzag,segment length=4pt,amplitude=1.2pt}]
    (1.05,0)--(3.0,0);
  \node[below] at (-1.05,-0.05) {$-E_M(\vec k)$};
  \node[below] at (1.05,-0.05) {$+E_M(\vec k)$};
  \draw[orange!85!black,very thick,->]
    (1.65,1.05) .. controls (2.2,0.82) and (2.2,0.28) .. (1.65,0.12);
  \draw[orange!85!black,very thick,->]
    (-1.65,-1.05) .. controls (-2.2,-0.82) and (-2.2,-0.28) .. (-1.65,-0.12);
  \node[orange!85!black,align=center] at (2.25,1.34)
    {positive-energy\\boundary value};
  \node[orange!85!black,align=center] at (-2.25,-1.34)
    {negative-energy\\boundary value};
\end{tikzpicture}
\end{center}

The output of the continuation is a family of Lorentzian time-ordered Green
functions with specified singularity prescriptions. The spectral poles and
cuts identify one-particle and continuum invariant masses seen by the field.
The next construction concerns asymptotic states: it asks for Hilbert-space
limits in which separated wave packets behave as freely propagating stable
particles in the far past or far future. That construction is logically
additional to the analytic continuation of Green functions.
